\section{True principal colength of every original finite torsion-free curve module}

\begin{definition}[The whole original basis principal quotient]
\label{def:fc-basis-principal-quotient}
\leanok
\lean{EnvelopeProperty.PrincipalModuleColength.basisPrincipalQuotientIso}
An actual original finite basis identifies the whole quotient by
multiplication by an original scalar with the product of the
original principal ring quotients. The isomorphism is constructed
from the original basis coordinate maps and original quotient maps.
\end{definition}
\begin{proof}
The actual basis maps the original scalar image onto the product
of the same original principal ideals. Original quotient descent
and the genuine quotient-of-products isomorphism give the whole
comparison, preserving the same scalar and every original basis
coordinate.
\end{proof}

\begin{theorem}[True principal colength on an original finite basis]
\label{thm:fc-basis-principal-colength}
\leanok
\lean{EnvelopeProperty.PrincipalModuleColength.length_quotient_scalar_of_basis}
\uses{def:fc-basis-principal-quotient}
The length of this SAME original principal module quotient is the
actual original basis cardinality times the true native ring order
of the SAME original scalar.
\end{theorem}
\begin{proof}
The preceding whole quotient isomorphism preserves the actual
module length. True finite-product length additivity computes the
length as the sum of the SAME original principal quotient lengths,
which by definition are the native original ring order.
\end{proof}

\begin{theorem}[Actual generic rank computes all original principal colengths]
\label{thm:fc-principal-module-colength-generic-rank}
\leanok
\lean{EnvelopeProperty.PrincipalModuleColength.length_quotient_scalar_eq_finrank_mul_ord}
\uses{thm:fc-basis-principal-colength,thm:fc-one-dimensional-full-free-lattice,thm:fc-finite-colength-endomorphism-index}
Over any original Noetherian domain of dimension at most one,
every original finite torsion-free module and every original
nonzero scalar satisfy: the true length of the principal module
quotient is the original generic module rank times the true native
ring order of that SAME scalar. No flatness, normality, regularity,
free basis or finite-colength input is assumed.
\end{theorem}
\begin{proof}
Internally produce the actual full generic-rank free submodule with
finite original colength. Original torsion-freeness and the
nonzero scalar make original multiplication injective. It preserves
this same original submodule. The proved finite-colength index
cancellation transfers the true quotient length to this internally
constructed free lattice; the preceding actual basis calculation
gives the original generic-rank multiple. This applies to nonflat
finite modules over singular original curve rings.
\end{proof}
